\documentclass[10pt,letterpaper]{article}
\usepackage{cornell-notes}

\title{Clause 05.02.12: Architecture Decisions And Rationale}
\author{}
\date{}
\setDocTitle{Clause 05.02.12: Architecture Decisions And Rationale}
\setDocSubtitle{ISO/IEC/IEEE 42010:2022}
\setDocOwner{ISO 42010 Cornell Notes}
\setDocDate{\today}
\setCornellCollection{ISO/IEC/IEEE 42010:2022 Cornell Notes}
\setCornellUnitType{Clause}
\setCornellUnitNumber{05.02.12}
\setCornellUnitTitle{Architecture Decisions And Rationale}
\hypersetup{pdftitle={Clause 05.02.12: Architecture Decisions And Rationale},pdfsubject={Cornell notes},pdfauthor={}}

\begin{document}
\maketitle
\begin{CornellOverviewBox}{Overview}
{\Large\bfseries\textcolor{CNavy}{Section 5.2.12: Architecture decisions and rationale}}\\[3pt]
{\small\textbf{Classification:} Conceptual foundation}\\
{\small\textbf{Learning objective:} Explain the role of architecture decisions and rationale and connect it to related architecture-description concepts.}
\end{CornellOverviewBox}
\vspace{3pt}

\begin{CornellNotesTable}
\CornellNoteRow{Purpose / key concept}{An architecture decision is a collection of choices made in the overall context of an architecture.}
\CornellNoteRow{Core details}{\begin{itemize}[leftmargin=*,nosep,topsep=1pt]
\item These choices usually pertain to various AD elements, entity requirements, or environmental influences on the architecture.
\item Architecture rationale records explanation, justification or reasoning about architecture decisions.
\item The rationale for a decision can include the following items: the basis for making a decision, impact on quality attributes, alternatives and trade-offs considered, potential consequences of the decision, architectural principles, and citations to sources of additional information.
\end{itemize}}
\CornellNoteRow{Distinction}{An architecture decision captures a choice; architecture rationale records the explanation, justification, and reasoning for that choice.}
\CornellNoteRow{Example / application}{Selection of architecture concepts, choice of AD elements, selection of ADFs, choice of architecture layering scheme, choice of underlying technology, choice of business components, choice of tactics to use for achieving system qualities, choice of business processes, choice of applicable patterns, choice of style(s) to be applied, choice of range of implementation technologies or other realizations to be considered, and choice of option sets to be considered.}
\CornellNoteRow{Interpretive note}{\begin{itemize}[leftmargin=*,nosep,topsep=1pt]
\item Requirements for capturing decisions and rationale within an AD are specified in 6.10.
\end{itemize}}
\CornellNoteRow{Related sections}{6.10}
\CornellNoteRow{Active recall}{\begin{itemize}[leftmargin=*,nosep,topsep=1pt]
\item What is the central role of Architecture decisions and rationale?
\item How does Architecture decisions and rationale connect to adjacent architecture-description concepts?
\item What closely related concepts must be kept distinct?
\end{itemize}}
\end{CornellNotesTable}


\begin{CornellSummaryBox}{Summary}
An architecture decision is a collection of choices made in the overall context of an architecture. An architecture decision captures a choice; architecture rationale records the explanation, justification, and reasoning for that choice.
\end{CornellSummaryBox}

\end{document}
